% ==================== PARTIE B — CONDITIONS ====================
% Étape de conception : règles de sécurité et choix du mode de charge.
%
% Progression :
%   B0  QCM — lire un if/else simple (sans piège)
%   B1  compléter UNE condition composée  (trou dans le code, escapeinside)
%   B2  repérer (numéroter) et corriger des erreurs
%   (QCM court-circuit déplacé en partie F ; conception supprimée, doublon de D)

\element{cond-lecture-if}{
  La borne surveille la température de son module de puissance avec le code suivant.

  \begin{question}{cond-lecture-if}
    \begin{minipage}{.45\textwidth}
      \lstinputlisting[language=c]{sources/codes/b0-lecture-if.c}
    \end{minipage}\hfill
    \begin{minipage}{.5\textwidth}
      Qu'affiche ce fragment de code ?
      \begin{reponses}
        \bonne{\texttt{Module OK}}
        \mauvaise{\texttt{Surchauffe : charge coupee}}
        \mauvaise{\texttt{Module OK} puis, sur la ligne suivante, \texttt{Surchauffe : charge coupee}}
        \mauvaise{Rien : aucune des deux branches n'est exécutée}
      \end{reponses}
    \end{minipage}
  \end{question}
}

\element{cond-completion}{
  \begin{question}{cond-completion}
    La borne n'autorise une charge que si \textbf{toutes} les conditions d'exploitation
    sont réunies :
  \begin{itemize}\setlength{\itemsep}{0pt}
      \item tension réseau comprise entre 360 V et 440 V, \textbf{bornes incluses} ;
      \item courant demandé \textbf{au plus} égal à 63 A ;
      \item température du module de puissance \textbf{strictement} inférieure à 65 °C.
    \end{itemize}

    Écrire la condition manquante \textbf{directement dans le cadre pointillé}.

    \lstinputlisting[language=c,escapeinside={(*}{*)}]{sources/codes/b2-autorisation-a-completer.c}

    Justifier brièvement le choix des opérateurs aux bornes de la plage de tension.

    \evaluationProfGrille{2}{3}
  \end{question}
}

\element{cond-correction}{
  \begin{question}{cond-correction}
    Un technicien a écrit le module de surveillance ci-dessous. Compilé sans option
    d'avertissement, il affiche parfois « Charge interrompue » alors qu'aucun défaut
    n'a été signalé.

    \lstinputlisting[language=c]{sources/codes/b3-securite-a-corriger.c}

    Repérer la ou les erreurs présentes dans ce programme en les \textbf{numérotant
    directement sur le listing} (1, 2\ldots). Puis, dans le cadre
    ci-dessous, donner pour chaque numéro la \textbf{ligne corrigée} et le comportement
    fautif que la correction supprime.

    \evaluationProfGrille{3}{7}
  \end{question}
}
